% !TEX root = ../main.tex
\documentclass[../Main.tex]{subfiles}

\begin{document}
\chapter*{Forord}
Fysikk er vanskelig. Noen ganger skikkelig vanskelig. Men jeg tror likevel at de aller fleste fint kan klare å få til minimum en C og gjerne en B i et brukerkurs i ingeniørfysikk. Men det må gjøres rett. Dette kompendiet er ment å skulle hjelpe deg til nettopp det. Mine beste tips for å lykkes i kurset er:
\begin{itemize}
    \item Fokuser på sammenhenger og forståelse. Fysikk ER forståelse. Pugging holder ikke engang til E.
    \item Forbered deg til forelesningene, og vær aktiv. Da husker du ting mye bedre.
    \item Jobb med oppgavene! Å selv måtte anvende det du har lært i oppgaveløsning er helt essensielt for å spikre fysikkunnskapen. Det er også gjennom å løse oppgaver at du finner ut hva du ikke kan, og dermed hva du må jobbe mer med.
    \item Før gode notater! Bruk tegninger aktivt, vær nøye med enheter, og skriv ned mellomregninger. Dette hjelper deg til å tenke klart. Husk at det handler ikke om å greie å løse én spesifikk oppgave. Det handler om å forstå løsningsstrategien godt nok til at man kan anvende den i nye oppgaver også.
\end{itemize}
\section*{Forfattere av kompendiet}
Jeg har fått god hjelp av flinnke studenter til å skrive notatene mine over i tex-format, og å utfylle de der det har trengtes. Tusen takk til Honja Kareem,  Lillian Ellefsen Valla og Mathias Vestgøte!

\section*{Foreslått framdriftsplan}
Dette kompendiet egner seg best til å brukes i kombinasjon med forelesningene og oppgavene som blir gitt i kurset. Jeg har laget en foreslått framdriftsplan for hvordan man kan bruke kompendiet sammen med forelesningene og oppgavene.
\begin{center}
\begin{tabular}{ll}
\hline
\textbf{Uke} & \textbf{Tema} \\
\hline
Uke 34 & \hyperref[Chapter:hvaerfysikk]{Kapittel 1: Hva er fysikk?} \\
Uke 35 & \hyperref[Chapter:matematikk]{Kapittel 2: Matematiske verktøy vi trenger} \\
Uke 36–37 & \hyperref[Chapter:kinematikk]{Kapittel 3: Kinematikk} \\
Uke 38–39 & \hyperref[Chapter:dynamikk]{Kapittel 4: Dynamikk} \\
Uke 40 & \hyperref[Chapter:energi]{Kapittel 5: Arbeid og energi} \\
Uke 41 & \hyperref[Chapter:statikk]{Kapittel 6: Statikk} \\
Uke 42–43 & \hyperref[Chapter:fluiddynamikk]{Kapittel 7: Fluiddynamikk} \\
Uke 44–45 & \hyperref[Chapter:termofysikk]{Kapittel 8: Termofysikk} \\
Uke 46–47 & \hyperref[Chapter:elektrisitet]{Kapittel 9: Elektrisitet} \\
\hline
\end{tabular}
\end{center}

\end{document}